\section{Task 1}
\makeatletter\def\@currentlabel{Task 1}\makeatother\label{sec:task1}

\subsection{Overview}
\begin{frame}{Overview}
  \begin{block}{Prompt}
    Classification of a dataset with 100 classes of different sports, on the clean test set.
  \end{block}
  \vspace{\baselineskip}
  The sheer number of classes and different variants of sports required us to make use of moderate augmnentation even in this task,
  to generate enough data samples to allow the model to learn correctly while trying to avoid overfitting and class confusion.
\end{frame}

\subsection{Model}
\begin{frame}{Model}
  \begin{itemize}
    \item Gestione valori mancanti
    \item Normalizzazione / encoding
    \item Train/val/test split
  \end{itemize}
\end{frame}
\subsection{Augmentation Techniques}
\begin{frame}{Augmentation Techniques}
  \begin{itemize}
    \item Aug1
    \item Aug2
    \item So on
  \end{itemize}
\end{frame}

\begin{frame}{Results}
\end{frame}
